\documentclass{article}
\usepackage{fontspec}
\usepackage{polyglossia}
\usepackage[unicodemath,all]{xamsmath}
\setmainfont{STIX Two Text}
\setmathfont{STIX Two Math}
\usepackage{graphicx}
\usepackage{xcolor}
\usepackage{hyperref}

% Load all plugins
\XAMSMathLoadPlugin{foundations}
\XAMSMathLoadPlugin{algebra}
\XAMSMathLoadPlugin{calculus}
\XAMSMathLoadPlugin{combinatorics}
\XAMSMathLoadPlugin{probability}

\DeclareBracedMathOperator*\Res{Res}


\title{xamsmath Package Demonstration (unicode-math)}
\author{xamsmath v0.1-alpha}
\date{}

\begin{document}

\newpage

\maketitle

\begin{abstract}
  \texttt{xamsmath} is a LaTeX package that extends \texttt{amsmath} with modern LaTeX3 features and enhanced mathematical notation. Main features include Prime/Index/Exponent (PIE) macros handling, automatic delimiter splitting, scalable operators, and a modular plugin system for different mathematical domains. 
\end{abstract}

\tableofcontents

\section{Core Package Features}

The xamsmath package provides several enhancements to mathematical typesetting. The PIE (Prime, Index, Exponent) based on \texttt{mathcommand} package system allows mathematical commands to properly absorb primes, subscripts, and superscripts that follow them. This provides more natural mathematical notation.

\begin{verbatim}
% Two methods for PIE handling:
\newmathcommandPIE{\fallf}[4]  % PIE absorbed as last arguments
\piefication{\mycmd}           % PIE absorbed before regular arguments
\end{verbatim}

For complete details on PIE absorption, refer to the \texttt{mathcommand} package documentation.

The \verb|\DeclarePairedSplitDelimiter| command creates delimiters that automatically split at vertical bars, ideal for set notation and conditional expressions.

\begin{verbatim}
\DeclarePairedSplitDelimiter{\set}{\{}{\}}
$\set{x | x > 0}$                    % Automatic splitting
$\set[\Big]{x | x^2 > 1}$            % Explicit size  
$\set*{\frac{1}{x} | x \neq 0}$      % Auto-scaling
\end{verbatim}

The starred version \verb|\set*| disables auto-scaling delimiters while the unstarred version uses fixed sizing. This behavior can be configured via package options.

\begin{align*}
  \set*{\frac{1}{x} | x \neq 0} &\quad\text{No auto-scaling} \\
  \set[\Big]{x | x^2 > 1} 		&\quad\text{Explicit sizing} \\
  \set{\frac{1}{x} | x^2 > 0} 	&\quad\text{Auto-scaling} 
\end{align*}

% Large operators can be dynamically scaled using the scaling package option. Starred versions provide scaled operators while unstarred versions maintain normal sizing.
%\begin{verbatim}
%$\sum_{i=1}^n x_i$                    % Normal size
%$\sum*_{i = 1}^{n}{\frac{1}{x_i^2}}$              % Scaled to match fraction
%\end{verbatim}

%\begin{align*}
%  \sum_{i = 1}^{n}{\frac{\displaystyle{{\int}\frac{\dif x_i}{x_i^2}}}{x_i}} &\quad\text{Normal sum operator} \\
%  \sum*_{i = 1}^{n}{\frac{\displaystyle{{\int}\frac{\dif x_i}{x_i^2}}}{x_i}} &\quad\text{Scaled sum operator}
%\end{align*}

The \verb|\DeclareBracedMathOperator| command creates operators that can handle different brace types with PIE support and automatic sizing.

\begin{verbatim}
\DeclareBracedMathOperator{\Res}{Res}
$\Res_{z = 0} \int f(z)$                         % No braces
$\Res_{z=0}[f(z)]$                   % Square brackets
$\Res_{z=0}(f(z))$                   % Parentheses
$\Res_{z=0}\{f(z)\}$                 % Braces
\end{verbatim}

\begin{align*}
  \Res_{\xi = 0} \int f(\xi) &\quad\text{No braces} \\
  \Res_{\xi=0}[\int f(\xi)] &\quad\text{With square brackets} \\
  \Res_{\xi=0}(\int f(\xi)) &\quad\text{With parentheses} \\
  \Res_{\xi=0}{\int f(\xi)} &\quad\text{With braces}
\end{align*}

\section{Plugins}

Each plugin is loaded with \verb|\XAMSMathLoadPlugin{<name>}|. The examples below show the source code as it appears in a document, followed immediately by the rendered output.

\begin{verbatim}
\XAMSMathLoadPlugin{foundations}
\XAMSMathLoadPlugin{algebra}
\XAMSMathLoadPlugin{calculus}
\XAMSMathLoadPlugin{combinatorics}
\XAMSMathLoadPlugin{probability}
\end{verbatim}

\subsection{Foundations Plugin}

The foundations plugin provides notation for logic, set theory, mappings, and topology.

\subsubsection*{Logic}
\begin{verbatim}
\begin{align*}
p \limp q &\quad\text{Implication} &
p \liff q &\quad\text{Equivalence} \\
\nec p &\quad\text{Necessity} &
\poss p &\quad\text{Possibility}
\end{align*}
\end{verbatim}
\begin{align*}
p \limp q &\quad\text{Implication} &
p \liff q &\quad\text{Equivalence} \\
\nec p &\quad\text{Necessity} &
\poss p &\quad\text{Possibility}
\end{align*}

\subsubsection*{Set theory and mappings}
\begin{verbatim}
\begin{align*}
\set{x | x > 0} &\quad\text{Set notation} &
\card{A} &\quad\text{Cardinality} \\
\powerset{X} &\quad\text{Power set} &
A \symdiff B &\quad\text{Symmetric difference} \\
f: X \injto Y &\quad\text{Injective function} &
\complement{A} &\quad\text{Set complement}
\end{align*}
\end{verbatim}
\begin{align*}
\set{x | x > 0} &\quad\text{Set notation} &
\card{A} &\quad\text{Cardinality} \\
\powerset{X} &\quad\text{Power set} &
A \symdiff B &\quad\text{Symmetric difference} \\
f: X \injto Y &\quad\text{Injective function} &
\complement{A} &\quad\text{Set complement}
\end{align*}

\subsubsection*{Topology}
\begin{verbatim}
\begin{align*}
\Cl{A} &\quad\text{Closure} &
\Int{A} &\quad\text{Interior} \\
\Fr{A} &\quad\text{Boundary} &
\Hom(X) &\quad\text{Homology groups}
\end{align*}
\end{verbatim}
\begin{align*}
\Cl{A} &\quad\text{Closure} &
\Int{A} &\quad\text{Interior} \\
\Fr{A} &\quad\text{Boundary} &
\Hom(X) &\quad\text{Homology groups}
\end{align*}

\subsection{Algebra Plugin}

The algebra plugin covers complex numbers, number systems, linear algebra, vectors, and number-theoretic delimiters.

\subsubsection*{Complex numbers}
\begin{verbatim}
\begin{align*}
\conj{z} &\quad\text{Conjugate} &
\abs{z} &\quad\text{Modulus} \\
\Re(z) &\quad\text{Real part} &
\Im(z) &\quad\text{Imaginary part}
\end{align*}
\end{verbatim}
\begin{align*}
\conj{z} &\quad\text{Conjugate} &
\abs{z} &\quad\text{Modulus} \\
\Re(z) &\quad\text{Real part} &
\Im(z) &\quad\text{Imaginary part}
\end{align*}

\subsubsection*{Number systems}
\begin{verbatim}
\[
\N, \Z, \Q, \R, \C, \H
\]
\end{verbatim}
\[
\N, \Z, \Q, \R, \C, \H
\]

\subsubsection*{Linear algebra}
\begin{verbatim}
\begin{align*}
\norm{v} &\quad\text{Vector norm} &
\norm[\Big]{Av} &\quad\text{Explicitly sized norm} \\
\spr{u}{v} &\quad\text{Scalar product} &
\cpr{u}{v} &\quad\text{Cross product} \\
\tr(A) &\quad\text{Trace} &
\rk(A) &\quad\text{Rank}
\end{align*}
\end{verbatim}
\begin{align*}
\norm{v} &\quad\text{Vector norm} &
\norm[\Big]{Av} &\quad\text{Explicitly sized norm} \\
\spr{u}{v} &\quad\text{Scalar product} &
\cpr{u}{v} &\quad\text{Cross product} \\
\tr(A) &\quad\text{Trace} &
\rk(A) &\quad\text{Rank}
\end{align*}

\subsubsection*{Vector shorthand}
\begin{verbatim}
\[
\va, \vb, \vc, \vi, \vj, \vx, \vy, \vz
\]
\end{verbatim}
\[
\va, \vb, \vc, \vi, \vj, \vx, \vy, \vz
\]

\subsubsection*{Number theory}
\begin{verbatim}
\begin{align*}
\ceil{x} &\quad\text{Ceiling function} &
\floor{x} &\quad\text{Floor function} \\
\ceil*[\Big]{\frac{x}{2}} &\quad\text{Explicitly sized ceiling} &
a \divby b &\quad\text{Divisibility}
\end{align*}
\end{verbatim}
\begin{align*}
\ceil{x} &\quad\text{Ceiling function} &
\floor{x} &\quad\text{Floor function} \\
\ceil*[\Big]{\frac{x}{2}} &\quad\text{Explicitly sized ceiling} &
a \divby b &\quad\text{Divisibility}
\end{align*}

\subsection{Calculus Plugin}

The calculus plugin provides derivative commands, differential symbols, interval notation, and vector-calculus operators.

\subsubsection*{Derivatives}
\begin{verbatim}
\begin{align*}
\od{y}{x} &\quad\text{Ordinary derivative} &
\od^2_{x=0}{y}{x} &\quad\text{Second derivative at a point} \\
\pd{u}{t} &\quad\text{Partial derivative} &
\pd^2{u}{x \partial y} &\quad\text{Mixed partial derivative}
\end{align*}
\end{verbatim}
\begin{align*}
\od{y}{x} &\quad\text{Ordinary derivative} &
\od^2_{x=0}{y}{x} &\quad\text{Second derivative at a point} \\
\pd{u}{t} &\quad\text{Partial derivative} &
\pd^2{u}{x \partial y} &\quad\text{Mixed partial derivative}
\end{align*}

\subsubsection*{Derivative fraction styles}
\begin{verbatim}
\begin{align*}
\lod{y}{x} &\quad\text{Small fraction derivative} &
\tod{y}{x} &\quad\text{Text-style fraction} \\
\lpd{u}{t} &\quad\text{Small fraction partial} &
\tpd{u}{t} &\quad\text{Text-style partial}
\end{align*}
\end{verbatim}
\begin{align*}
\lod{y}{x} &\quad\text{Small fraction derivative} &
\tod{y}{x} &\quad\text{Text-style fraction} \\
\lpd{u}{t} &\quad\text{Small fraction partial} &
\tpd{u}{t} &\quad\text{Text-style partial}
\end{align*}

\subsubsection*{Multivariable calculus}
\begin{verbatim}
\begin{align*}
\grad f &\quad\text{Gradient} &
\rot \vec{F} &\quad\text{Curl} \\
\div \vec{F} &\quad\text{Divergence} &
\interval{a}{b} &\quad\text{Open interval}
\end{align*}
\end{verbatim}
\begin{align*}
\grad f &\quad\text{Gradient} &
\rot \vec{F} &\quad\text{Curl} \\
\div \vec{F} &\quad\text{Divergence} &
\interval{a}{b} &\quad\text{Open interval}
\end{align*}

\subsection{Combinatorics Plugin}

The combinatorics plugin provides binomial variants, Stirling numbers, and PIE-aware falling and rising factorials.

\subsubsection*{Combinatorial numbers}
\begin{verbatim}
\begin{align*}
\multbinom{n}{k} &\quad\text{Multiset coefficient} \\
\Stirling{n}{k} &\quad\text{Second-kind Stirling number} \\
\stirling{n}{k} &\quad\text{First-kind Stirling number}
\end{align*}
\end{verbatim}
\begin{align*}
\multbinom{n}{k} &\quad\text{Multiset coefficient} \\
\Stirling{n}{k} &\quad\text{Second-kind Stirling number} \\
\stirling{n}{k} &\quad\text{First-kind Stirling number}
\end{align*}

\subsubsection*{Factorials}
\begin{verbatim}
\begin{align*}
\fallf{x}^{n} &\quad\text{Falling factorial} &
\fallf{x}_{k}^{n} &\quad\text{Falling factorial with index} \\
\risef{x}^{n} &\quad\text{Rising factorial} &
\risef{x}^{n}_{k} &\quad\text{Rising factorial with index}
\end{align*}
\end{verbatim}
\begin{align*}
\fallf{x}^{n} &\quad\text{Falling factorial} &
\fallf{x}_{k}^{n} &\quad\text{Falling factorial with index} \\
\risef{x}^{n} &\quad\text{Rising factorial} &
\risef{x}^{n}_{k} &\quad\text{Rising factorial with index}
\end{align*}

\subsection{Probability Plugin}

The probability plugin provides conditional probability, expectation-like operators, covariance, distributions, and statistical shorthand.

\subsubsection*{Probability measures}
\begin{verbatim}
\begin{align*}
\P(A) &\quad\text{Probability measure} &
\P(A | B) &\quad\text{Conditional probability} \\
\P_{X}(A) &\quad\text{Under distribution X} &
\Law(X) &\quad\text{Law of a random variable}
\end{align*}
\end{verbatim}
\begin{align*}
\P(A) &\quad\text{Probability measure} &
\P(A | B) &\quad\text{Conditional probability} \\
\P_{X}(A) &\quad\text{Under distribution X} &
\Law(X) &\quad\text{Law of a random variable}
\end{align*}

\subsubsection*{Expectation and moments}
\begin{verbatim}
\begin{align*}
\E[X] &\quad\text{Expectation} &
\E_P^2[X | Y] &\quad\text{Conditional expectation} \\
\E_{\P}[X] &\quad\text{Under measure P} &
\D[X] &\quad\text{Variance}
\end{align*}
\end{verbatim}
\begin{align*}
\E[X] &\quad\text{Expectation} &
\E_P^2[X | Y] &\quad\text{Conditional expectation} \\
\E_{\P}[X] &\quad\text{Under measure P} &
\D[X] &\quad\text{Variance}
\end{align*}

\subsubsection*{Statistics and distributions}
\begin{verbatim}
\begin{align*}
\Cov{X}{Y} &\quad\text{Covariance} &
\std[X] &\quad\text{Standard deviation} \\
X \sim \Normal(\mu,\sigma^2) &\quad\text{Normal distribution} &
X \sim \Uniform(a,b) &\quad\text{Uniform distribution}
\end{align*}
\end{verbatim}
\begin{align*}
\Cov{X}{Y} &\quad\text{Covariance} &
\std[X] &\quad\text{Standard deviation} \\
X \sim \Normal(\mu,\sigma^2) &\quad\text{Normal distribution} &
X \sim \Uniform(a,b) &\quad\text{Uniform distribution}
\end{align*}

\subsection{Advanced Examples}

The plugins can be composed in larger expressions.

\begin{verbatim}
\begin{align*}
\od^2_{\xi=0}{y}{\xi} + \pd{u}{t} &= \grad^2 u \\
\sum_{k=0}^n \Stirling{n}{k} \risef{\xi}^k &= \xi^n \\
\E[\Cov{\xi}{\eta} | \zeta] &= 0 \\
\set{f: X \bijto Y | \card{X} = \card{Y}} &\quad\text{Bijections}
\end{align*}
\end{verbatim}
\begin{align*}
\od^2_{\xi=0}{y}{\xi} + \pd{u}{t} &= \grad^2 u \\
\sum_{k=0}^n \Stirling{n}{k} \risef{\xi}^k &= \xi^n \\
\E[\Cov{\xi}{\eta} | \zeta] &= 0 \\
\set{f: X \bijto Y | \card{X} = \card{Y}} &\quad\text{Bijections}
\end{align*}

\section{Package Configuration}

xamsmath supports several options to enable specific features:

\begin{verbatim}
\usepackage[unicodemath]{xamsmath}
\usepackage[fonts]{xamsmath}
\usepackage[symbols]{xamsmath}
\usepackage[operators]{xamsmath}
\usepackage[braces]{xamsmath}
\usepackage[eqmicrotype]{xamsmath}
\usepackage[all]{xamsmath}
\end{verbatim}

The brace handling option swaps the behavior of starred and unstarred commands in paired delimiters, making unstarred commands use auto-scaling by default.

\end{document}
